Le tambour d'une machine à laver est entraîné par un moteur électrique. La
transmission du mouvement est assurée par une courroie tournant sans
glissement. La fréquence de rotation du moteur est
$f_{A}=\qty{3000}{tr\per\min}$.

La poulie du moteur a un diamètre $D_{A}=\qty{10}{\cm}$ et celui de la poulie du
tambour est $D_{B}=\qty{40}{\cm}$.
\begin{figure}[H]
    \centering
    \begin{tikzpicture}[>=Stealth]
        \newdimen\ri \ri=4mm
        \newdimen\rii \rii=8mm
        \newdimen\R \R=15mm
        \node[circle,draw,fill=lightgray,inner sep=0pt,
            minimum size=2\R] (S) {};
        \foreach \x/\r/\l [count=\i] in {-6/\ri/A,0/\rii/B}
        \node[circle,draw,fill=gray,inner sep=0pt,minimum size=2\r,
            label=below:poulie~\l] (P\i) at (\x,0) {};
        \draw[thick] (P2.95) arc (95:-95:\rii) --
            (P1.-95) arc (-95:-270+5:\ri) --
            node[above,sloped,pos=0.4] {courroie} (P2.north);
        \foreach \i in {1,2}
            \node[circle,fill,inner sep=1pt] at (P\i.center) {};
        \node[above=1mm] at (P1.north) {moteur};
        \node[above=1mm] at (S.north) {tambour};
    \end{tikzpicture}
    \caption{}
    \label{fig:scie}
\end{figure}
\begin{enumerate}
    \item Calculer la fréquence de rotation du moteur en tours par seconde.
    \item Déterminer la vitesse angulaire $\omega_{A}$ du moteur en $\unit{\w}$.
    \item Calculer la vitesse linéaire d'un point de la courroie en
          $\unit{\v}$.
    \item Déterminer la vitesse angulaire $\omega_{B}$ du tambour.
    \item En déduire la relation littérale entre les fréquences de rotation
          $f_{A}$ et $f_{B}$ du moteur et du tambour. Calculer $f_{B}$ en
          $\unit{\tr\per\min}$.
    \item Calculer la vitesse d'un point de la circonférence du tambour de
          diamètre $D_{T}=\qty{100}{\cm}$.
\end{enumerate}
